\section{函数项级数和幂级数的逐项求导}
\label{sec:ru-termwise-differentiation}

有限和的导数等于各项导数之和。对函数项级数，能否同样逐项求导，就是能否交换“求导”与“取极限”两个操作。
上一章已经看到一致收敛能保持连续性，但下面的例子说明，它不足以保持导数。

令$f_n(x)=\sin(n^2x)/n$。由$|f_n(x)|\leq1/n$，$f_n$在$\mathbb R$上一致收敛到$0$，极限函数的导数为$0$；但$f_n'(x)=n\cos(n^2x)$，在$x=0$处等于$n\to\infty$，并不收敛。
函数值的振幅虽小，振荡的频率却增长得更快，斜率因而越来越大。要交换求导与极限，需要直接控制导数。
\begin{theorem}[函数列的求导与极限交换]
\label{thm:ru-uniform-derivative-limit}
设$f_n$在$[a,b]$上连续可微，导函数列$(f_n')$在$[a,b]$上一致收敛到$g$，并且存在一点$x_0\in[a,b]$使$(f_n(x_0))$收敛。则$f_n$在$[a,b]$上一致收敛到某个可微函数$f$，并且$f'=g$。
\end{theorem}
\begin{proof}
先证$(f_n)$一致收敛。对函数$f_n-f_m$用中值定理，
\[
 |f_n(x)-f_m(x)|\leq|f_n(x_0)-f_m(x_0)|+\sup_{[a,b]}|f_n'-f_m'|\cdot(b-a),
\]
右侧与$x$无关，并且由两个条件可以任意小，由一致柯西准则，$f_n$一致收敛到某个$f$。

再证$f'(c)=g(c)$。对$x\neq c$，记$Q_n(x)=\frac{f_n(x)-f_n(c)}{x-c}$，$Q(x)=\frac{f(x)-f(c)}{x-c}$。对$f_n-f_m$在$c$与$x$之间用中值定理，得$|Q_n(x)-Q_m(x)|\leq\sup|f_n'-f_m'|$，
令$m\to\infty$，由$f_m'$一致收敛到$g$，得$|Q_n(x)-Q(x)|\leq\sup_{[a,b]}|f_n'-g|$，右侧与$x$无关且趋于零，即$Q_n\to Q$对$x$一致。于是
\[
 |Q(x)-g(c)|\leq|Q(x)-Q_n(x)|+|Q_n(x)-f_n'(c)|+|f_n'(c)-g(c)|.
\]
给定$\varepsilon$，先取$n$使第一项和第三项对所有$x$都小于$\varepsilon/3$，再由$f_n'(c)$的定义取$x$足够接近$c$使第二项小于$\varepsilon/3$，所以$Q(x)\to g(c)$。
\end{proof}
证明与上一章连续性保持的论证结构相同：把总误差拆成三段，关键是用一致收敛让同一个$n$对所有$x$有效。这里被一致控制的是差商，它由导数的一致收敛通过中值定理得到。
一点处的收敛条件不能省去：$f_n\equiv n$的导数恒为零，一致收敛到$0$，函数列本身却发散。它排除的正是整列函数一起平移、导数却不变的情形。

把定理用于部分和$S_N=\sum_{k\leq N}g_k$：若各项$g_k$在$[a,b]$上连续可微，导数级数$\sum g_k'$在$[a,b]$上一致收敛，且$\sum g_k(x_0)$在某一点收敛，
则$\sum g_k$在$[a,b]$上一致收敛，和函数可微，并且可以逐项求导：$\bigl(\sum g_k\bigr)'=\sum g_k'$。两个条件各有分工：导数级数的一致收敛控制差商，一点处的收敛确定和函数的取值。

幂级数满足这些条件。设$\sum c_k(x-x_0)^k$的收敛半径为$R>0$，逐项求导得到$\sum kc_k(x-x_0)^{k-1}$。由$e^u\geq u^2/2$（取$u=\log k$）得$0\leq(\log k)/k\leq2/\log k\to0$，从而$k^{1/k}=e^{(\log k)/k}\to1$，于是
\[
 \limsup_{k\to\infty}|kc_k|^{1/k}=\limsup_{k\to\infty}|c_k|^{1/k},
\]
由Cauchy--Hadamard公式，导数级数与原级数有相同的收敛半径（整体乘以$x-x_0$不改变收敛半径）。导数级数也是幂级数，在收敛区间的每个内部闭子区间上一致收敛，原级数在中心收敛，于是在这些闭子区间上可以逐项求导。
每一点都落在某个这样的闭子区间中，所以和函数$S$在整个收敛区间内可微，且$S'(x)=\sum_kkc_k(x-x_0)^{k-1}$。

导数级数仍是收敛半径为$R$的幂级数，可以继续逐项求导。于是和函数在收敛区间内无穷次可微，代入$x=x_0$得$S^{(k)}(x_0)=k!\,c_k$：幂级数就是它的和函数的Taylor级数。
例如对$\sum x^k=1/(1-x)$逐项求导，得$\sum_{k\geq1}kx^{k-1}=1/(1-x)^2$（$|x|<1$）；对指数级数逐项求导，又得到$(e^x)'=e^x$，与第~\ref{sec:ru-elementary-derivatives}节由基本极限得到的结果一致。
端点处的求导不能从区间内部的结论自动得到：$\sum x^k/k^2$在$[-1,1]$上收敛，逐项求导得到的$\sum x^{k-1}/k$却在$x=1$处发散。
